\section{Classes Modulo \texorpdfstring{$m$}{m}}

\begin{proposition}
\label{prop:euclidean-division-unique-remainder}
$\forall a \in \mathbb{Z}$ (the dividend) and $\forall m \in \mathbb{N}$ so that $m > 1$ (the modulus), $\exists! q, r \in \mathbb{Z} : a = mq + r$, with $0 \leq r < m$.
\end{proposition}

\begin{center}
    \line(1,0){450}
\end{center}

\begin{theorem}[Existence of a Unique Remainder]
$\forall a \in \mathbb{Z}$ and $\forall m \in \mathbb{N} : m > 1$, $\exists! r \in \{0, 1, \dots, m-1\} : a \equiv r \pmod{m}$.
\end{theorem}
\begin{proof}
By Proposition \ref{prop:euclidean-division-unique-remainder} (Euclidean Division), we have $\exists! q, r \in \mathbb{Z} : a = mq + r$, where $0 \leq r < m$.
Rearranging this equation gives $a - r = mq$.
By the formal definition of divisibility, $m \mid (a - r)$.
By the formal definition of congruence, this directly implies $a \equiv r \pmod{m}$. The uniqueness of $r$ is guaranteed by the Euclidean Division Theorem. \qedhere
\end{proof}

\begin{center}
    \line(1,0){450}
\end{center}

\begin{theorem}[Congruence is an Equivalence Relation]
The relation of congruence modulo $m$ (denoted by $\equiv \pmod{m}$) is an equivalence relation on the set of integers $\mathbb{Z}$.
\end{theorem}
\begin{proof}
We must demonstrate that congruence modulo $m$ satisfies reflexivity, symmetry, and transitivity.
\begin{itemize}
    \item \textbf{Reflexivity:} $\forall a \in \mathbb{Z}$, $a - a = 0$. Since $m \mid 0$, $a \equiv a \pmod{m}$. The relation is reflexive.
    \item \textbf{Symmetry:} Assume $a \equiv b \pmod{m}$. $\implies m \mid (a - b)$, so $\exists k \in \mathbb{Z} : a - b = k m$.
    Multiplying by $-1$ yields $b - a = (-k) m$. Since $-k \in \mathbb{Z}$, $\implies m \mid (b - a)$, so $b \equiv a \pmod{m}$. The relation is symmetric.
    \item \textbf{Transitivity:} Assume $a \equiv b \pmod{m}$ and $b \equiv c \pmod{m}$.
    $\implies \exists k_1, k_2 \in \mathbb{Z} : a - b = k_1 m \text{ and } b - c = k_2 m$.
    Adding the equations:
    $$ (a - b) + (b - c) = k_1 m + k_2 m $$
    $$ a - c = (k_1 + k_2) m $$
    Since $k_1 + k_2 \in \mathbb{Z}$, $\implies m \mid (a - c)$, so $a \equiv c \pmod{m}$. The relation is transitive. \qedhere
\end{itemize}
\end{proof}

\begin{center}
    \line(1,0){450}
\end{center}

\begin{proposition}[Equivalence Classes Form a Partition]
Let $\sim$ be an equivalence relation on a set $X$. The set of equivalence classes, $X/\sim$, forms a partition of $X$.
\end{proposition}

\begin{center}
    \line(1,0){450}
\end{center}